\chapter{Foundations: lattices, completely monotone functions and formal verification}
\chaptermark{Foundations}
\label{ch1:chap}

\begin{quote}\small\itshape
This chapter sets the stage. It assumes only first-year undergraduate analysis and linear algebra,
plus the basic theory of Fourier series. The proofs are written by hand, as one would write them in an exam;
the machine (Lean~4) enters only in the last part, to say what ``computer-verified'' means and, above all,
what it does not mean.
\end{quote}

\section*{How to read this book}

Every claim carries a \emph{tier} label that says \emph{who checked what}.
These tiers, taken from the Elenchus method~\cite{elenchus} and detailed in §\ref{ch1:sec:tiers}, are the following.
\begin{itemize}
\item \tierA{}: a statement checked by the Lean~4 kernel, with three controls: axioms on the whitelist, kernel replay (\lean{leanchecker}) and a search for backdoors (§\ref{ch1:sec:lean}).
\item \tierAs{}: as \tierA{}, but without the kernel replay (§\ref{ch1:sec:lean}).
\item \tierL{}: a result from the literature, read and precisely referenced.
\item \tierC{}: a human argument, written out in detail in this book, but not machine-checked.
\item \tierX{}: a floating-point numerical exploration; it illustrates, it never proves.
\end{itemize}

\begin{warnbox}[: honesty about where the results come from]
All results in this book were produced by AI agents (Claude) under the author's direction, and then audited
by a \emph{separate} AI agent. \textbf{No human has audited these results.} Moreover, everything that relies on the
\texttt{openai/math} repository is \emph{conditional} on that repository's theorems, which are themselves AI-generated and have not been
peer reviewed~\cite{oai}. A \tierA{} result guarantees that the \emph{Lean} statement follows from the
allowed axioms; it does not guarantee that this statement faithfully expresses the mathematical or physical statement one has in mind.
\end{warnbox}

\section{Lattices in the plane}
\label{ch1:sec:reseaux}

\subsection{Lattice, covolume, density}

\begin{definition}[Lattice in the plane]\label{ch1:def:reseau}
Let $(u,v)$ be a basis of $\R^2$. The \emph{lattice} it generates is $\Lambda=\Z u\oplus\Z v=\{mu+nv : (m,n)\in\Z^2\}$.
Its \emph{covolume} is $\operatorname{covol}(\Lambda)=|\det(u,v)|$, the area of the parallelogram
$P=\{su+tv : s,t\in[0,1)\}$, called the \emph{fundamental cell}. The covolume does not depend on the chosen basis
(two bases differ by a matrix in $\mathrm{GL}_2(\Z)$, of determinant $\pm1$).
\end{definition}

We write $N_\Lambda(R)=\#\{x\in\Lambda : |x|\le R\}$ for the number of points of $\Lambda$ in the closed disk of radius $R$.
A lattice is an example of a \emph{locally finite configuration}: every disk contains finitely many points.

\begin{proposition}[Density of a lattice]\label{ch1:prop:comptage}
If $V=\operatorname{covol}(\Lambda)$, then $N_\Lambda(R)/(\pi R^2)\to 1/V$ as $R\to\infty$. More precisely, with
$\delta=|u|+|v|$ and for $R>\delta$,
\[
  \pi (R-\delta)^2\ \le\ V\,N_\Lambda(R)\ \le\ \pi (R+\delta)^2 .
\]
\tierC{}
\end{proposition}

\begin{humanproof}
The translates $x+P$, $x\in\Lambda$, \emph{tile} the plane: every point $y=su'+tv'$ can be written uniquely as
$y=x+p$ with $x\in\Lambda$ (integer part of the coordinates in the basis $(u,v)$) and $p\in P$. Each cell
$x+P$ has area $V$ and diameter at most $\delta$.
\begin{itemize}
\item If $|x|\le R$, then $x+P\subset \overline{D}(0,R+\delta)$. The cells are disjoint, hence
$V\,N_\Lambda(R)\le \pi(R+\delta)^2$.
\item If $|y|\le R-\delta$, the point $y$ lies in a cell $x+P$ with $|y-x|\le\delta$, hence $|x|\le R$.
Thus $\overline{D}(0,R-\delta)\subset\bigcup_{|x|\le R}(x+P)$ and $\pi(R-\delta)^2\le V\,N_\Lambda(R)$.
\end{itemize}
Divide by $\pi R^2$ and let $R$ tend to infinity.
\end{humanproof}

We say that $\Lambda$ has \emph{density} $1/V$. A lattice of covolume $1$ has density $1$: this is the normalisation
used throughout the book, which allows different lattices to be compared at an equal number of points per unit area.

\subsection{The triangular lattice}

\begin{definition}[Triangular lattice $A$]\label{ch1:def:A}
Let $a>0$ with $a^2=2/\sqrt3$. The \emph{triangular} (or hexagonal) \emph{lattice} of covolume $1$ is
\[
  A=\Z\,u\oplus\Z\,v,\qquad u=a(1,0),\quad v=a\bigl(\tfrac12,\tfrac{\sqrt3}{2}\bigr).
\]
\end{definition}

The choice of $a$ comes from the covolume: $\det(u,v)=a^2\sqrt3/2=1$. Numerically $a\approx 1.0746$.

\begin{proposition}[Norms in $A$]\label{ch1:prop:normesA}
For $(m,n)\in\Z^2$ we have $|mu+nv|^2=a^2\,(m^2+mn+n^2)$. In particular the minimal distance between two points of $A$ is $a$, attained
by exactly $6$ vectors (so the central point has $6$ nearest neighbours, at the vertices of a regular hexagon). \tierC{}
\end{proposition}

\begin{humanproof}
Since $|u|^2=|v|^2=a^2$ and $u\cdot v=a^2/2$, we have $|mu+nv|^2=a^2(m^2+n^2+mn)$. Now
$m^2+mn+n^2=(m+\tfrac n2)^2+\tfrac34n^2\ge0$, with equality only for $m=n=0$. The integer solutions of
$m^2+mn+n^2=1$ satisfy $\tfrac34 n^2\le 1$, hence $n\in\{-1,0,1\}$; solving each case gives $(\pm1,0)$, $(0,\pm1)$,
$\pm(1,-1)$. There are six of them, and the value $1$ is the smallest nonzero value because $m^2+mn+n^2$ is a nonnegative integer.
\end{humanproof}

The first values of $k=m^2+mn+n^2$ and their multiplicities, compared with those of the square lattice $\Z^2$ (also of covolume $1$),
are given in Table~\ref{ch1:tab:couches}. These numbers were counted by machine computation \tierX{}; they are classical.

\begin{table}[ht]
\centering
\begin{tabular}{@{}llll@{}}
\toprule
Lattice & $k$ (value of the form) & multiplicity & $|x|^2$ \\
\midrule
$A$ & $1,\ 3,\ 4,\ 7,\ 9$ & $6,\ 6,\ 6,\ 12,\ 6$ & $a^2k\approx 1.155,\ 3.464,\ 4.619,\ 8.083,\ 10.392$ \\
$\Z^2$ & $1,\ 2,\ 4,\ 5,\ 8$ & $4,\ 4,\ 4,\ 8,\ 4$ & $1,\ 2,\ 4,\ 5,\ 8$ \\
\bottomrule
\end{tabular}
\caption{First shells of the lattices $A$ and $\Z^2$ (covolume $1$).}
\label{ch1:tab:couches}
\end{table}

The triangular lattice has neighbours that are farther away ($|x|^2\approx1.155$ against $1$) and more numerous ($6$ against $4$):
it ``balances'' the repulsion better, and this is what the universal optimality theorem will quantify.

\subsection{Dual lattice}

\begin{definition}[Dual lattice]\label{ch1:def:dual}
If $\Lambda=M\Z^2$ with $M$ the matrix with columns $u,v$, the \emph{dual lattice} is
$\Lambda^*=\{y\in\R^2 : x\cdot y\in\Z\ \ \forall x\in\Lambda\}=M^{-T}\Z^2$. Its covolume is $1/\operatorname{covol}(\Lambda)$.
\end{definition}

\begin{lemma}[$A$ is similar to its dual]\label{ch1:lem:autodual}
The lattice $A^*$ is the image of $A$ under a rotation; in particular $A$ and $A^*$ have the same norms, with the same multiplicities. \tierC{}
\end{lemma}

\begin{humanproof}
Let $M=a\begin{pmatrix}1&1/2\\0&\sqrt3/2\end{pmatrix}$. A direct computation gives
$M^{-T}=\frac1a\begin{pmatrix}1&0\\-1/\sqrt3&2/\sqrt3\end{pmatrix}$. Its two columns $b_1,b_2$ satisfy
$|b_1|^2=|b_2|^2=\frac{1}{a^2}\cdot\frac43=\frac{\sqrt3}{2}\cdot\frac43=a^2$ and $b_1\cdot b_2=-\frac{1}{a^2}\cdot\frac23=-\frac{a^2}{2}$:
two vectors of norm $a$ at an angle of $120^\circ$. Then $(b_1,b_1+b_2)$ is a basis of $A^*$ consisting of two vectors
of norm $a$ at an angle of $60^\circ$, that is, the basis of $A$ rotated. Hence $A^*$ is the image of $A$ under a rotation.
\end{humanproof}

\subsection{Energy of a lattice}

We study identical particles that repel each other through a \emph{pair potential} $g$, a function of the \emph{square} $t=|x-y|^2$ of
the distance (the convention of~\cite{ckmrv} and~\cite{oai}: Gaussians $e^{-\alpha t}$ are then the basic building blocks).
The Riesz potential $|x-y|^{-s}$ reads $g(t)=t^{-s/2}$.

\begin{definition}[Energy of a lattice]\label{ch1:def:enerreseau}
For $g:(0,\infty)\to[0,\infty)$, the energy per particle of a lattice $\Lambda$ is
$E_g(\Lambda)=\sum_{x\in\Lambda\setminus\{0\}} g(|x|^2)\in[0,+\infty]$.
\end{definition}

For a locally finite configuration $C$ of density $1$, we define the \emph{lower energy} as the average, over the points
of the disk of radius $R$, of the interactions with the other points of the disk, followed by a limit inferior:
\[
  E_g(C;R)=\frac{1}{\#C_R}\sum_{x\in C_R}\ \sum_{y\in C_R\setminus\{x\}} g(|x-y|^2),\qquad E_g(C)=\liminf_{R\to\infty}E_g(C;R)\in[0,\infty],
\]
with $C_R=C\cap\overline D(0,R)$. Each pair is counted twice, which indeed gives $\sum_{a\neq0}g(|a|^2)$ for a lattice.
\tierC{} The exact definition used by the \texttt{openai/math} repository, in $[0,\infty]$ arithmetic, will be quoted in §\ref{ch1:sec:univ}.

\begin{proposition}[Riesz energy of a lattice]\label{ch1:prop:riesz}
For $g(t)=t^{-s/2}$ with $s>0$, the energy $E_g(\Lambda)=\sum_{x\ne0}|x|^{-s}$ is finite if and only if $s>2$. \tierC{}
\end{proposition}

\begin{humanproof}
Let $d>0$ be the smallest nonzero norm of $\Lambda$ and $V$ its covolume. There are finitely many nonzero points of norm $\le d$;
for $k\ge0$ let $S_k=\{x\in\Lambda : 2^kd<|x|\le 2^{k+1}d\}$. By Proposition~\ref{ch1:prop:comptage},
$\#S_k=N_\Lambda(2^{k+1}d)-N_\Lambda(2^kd)=\frac{3\pi d^2}{V}4^k+O(2^k)$, so there exist $0<c<C$ with $c\,4^k\le\#S_k\le C\,4^k$ for $k\ge k_0$.
On $S_k$ we have $(2^{k+1}d)^{-s}\le |x|^{-s}<(2^kd)^{-s}$. The sum over $S_k$ therefore lies between positive multiples of $4^k2^{-ks}=2^{k(2-s)}$.
The series $\sum_k 2^{k(2-s)}$ converges if and only if $s>2$.
\end{humanproof}

\begin{remark}
For $0<s\le2$ the energy of a lattice is infinite: the optimality inequality $\infty\le\infty$ is then true but void, and a meaningful
comparison would require a \emph{regularisation}. \tierC{}
\end{remark}

\subsection{Theta series and the Poisson summation formula}
\label{ch1:sec:poisson}

For $\alpha>0$, the \emph{Gaussian theta series} of the lattice $\Lambda$ is
\[
  \Theta_\Lambda(\alpha)=\sum_{x\in\Lambda}e^{-\pi\alpha|x|^2}\ \ (>1).
\]
(Not to be confused with the arithmetic Gauss sums $\sum e^{2\pi i n^2/p}$: here we are dealing with sums of \emph{real}
exponentials of quadratic forms.) The series converges, by comparison with a geometric series via the lower bound
$|mu+nv|^2\ge c(m^2+n^2)$. The energy of a Gaussian interaction $g(t)=e^{-\pi\alpha t}$ is therefore
$E_g(\Lambda)=\Theta_\Lambda(\alpha)-1$.

The key result is that $\Theta_\Lambda$ transforms simply when one passes from $\alpha$ to $1/\alpha$ and from $\Lambda$ to $\Lambda^*$.

\begin{theorem}[Poisson summation for Gaussians]\label{ch1:thm:poisson}
For every lattice $\Lambda$ in the plane of covolume $V$ and every $\alpha>0$,
\[
  \Theta_\Lambda(\alpha)=\frac{1}{V\,\alpha}\ \Theta_{\Lambda^*}\!\Bigl(\frac1\alpha\Bigr).
\]
For the triangular lattice ($V=1$, $A^*\simeq A$): $\Theta_A(\alpha)=\alpha^{-1}\Theta_A(1/\alpha)$. \tierC{}\
(The pointwise convergence of the Fourier series used in Step~3 is Fourier's theorem on the torus: \tierL{}.)
\end{theorem}

\begin{humanproof}[Human proof in three steps]
Convention: $\hat f(\xi)=\int f(x)e^{-2\pi i\,\xi\cdot x}\,dx$.

\emph{Step 1: Fourier transform of a Gaussian in dimension 1.} Let $\varphi(x)=e^{-\pi\alpha x^2}$ and $h=\hat\varphi$.
Differentiating under the integral sign (dominated by $|x|e^{-\pi\alpha x^2}$, which is integrable),
$h'(\xi)=\int(-2\pi i x)\varphi(x)e^{-2\pi i\xi x}dx$. Now $\varphi'(x)=-2\pi\alpha x\varphi(x)$, so $-2\pi ix\varphi(x)=\frac{i}{\alpha}\varphi'(x)$ and
\[
  h'(\xi)=\frac i\alpha\int\varphi'(x)e^{-2\pi i\xi x}dx=\frac i\alpha\Bigl([\varphi e^{-2\pi i\xi x}]_{-\infty}^{+\infty}+2\pi i\xi\int\varphi(x)e^{-2\pi i\xi x}dx\Bigr)=-\frac{2\pi\xi}{\alpha}\,h(\xi),
\]
(the bracket vanishes because $\varphi\to0$ at $\pm\infty$). This linear differential equation gives $h(\xi)=h(0)\,e^{-\pi\xi^2/\alpha}$, and
$h(0)=\int e^{-\pi\alpha x^2}dx=\alpha^{-1/2}$ (Gaussian integral $\int e^{-\pi x^2}dx=1$, computed in polar coordinates). Hence
$\widehat{e^{-\pi\alpha x^2}}(\xi)=\alpha^{-1/2}e^{-\pi\xi^2/\alpha}$. In dimension $2$, by Fubini,
$\widehat{e^{-\pi\alpha|x|^2}}(\xi)=\alpha^{-1}e^{-\pi|\xi|^2/\alpha}$.

\emph{Step 2: periodisation.} Write $\Lambda=M\Z^2$ and set $f(x)=e^{-\pi\alpha|Mx|^2}$ for $x\in\R^2$. The function
$F(x)=\sum_{n\in\Z^2}f(x+n)$ is well defined (locally normal convergence, together with that of its derivatives), $\mathcal C^\infty$, and $\Z^2$-periodic.
Its Fourier coefficients are, for $k\in\Z^2$,
\[
  c_k=\int_{[0,1]^2}F(x)e^{-2\pi ik\cdot x}dx=\int_{\R^2}f(x)e^{-2\pi ik\cdot x}dx=\hat f(k),
\]
(the translates of $[0,1]^2$ are put together to form $\R^2$). The change of variables $u=Mx$ ($dx=du/V$) gives
$\hat f(k)=\frac1V\int e^{-\pi\alpha|u|^2}e^{-2\pi i(M^{-T}k)\cdot u}du=\frac{1}{V\alpha}e^{-\pi|M^{-T}k|^2/\alpha}$ by Step~1.

\emph{Step 3: evaluation at $0$.} The $c_k$ are summable over $\Z^2$ and $F$ is continuous; the Fourier series of $F$ therefore converges normally
to a continuous function with the same coefficients, which is $F$ (uniqueness of the coefficients, Fejér's theorem). At $x=0$:
\[
  \sum_{n\in\Z^2}e^{-\pi\alpha|Mn|^2}=F(0)=\sum_{k\in\Z^2}c_k=\frac{1}{V\alpha}\sum_{k\in\Z^2}e^{-\pi|M^{-T}k|^2/\alpha}.
\]
The left-hand side is $\Theta_\Lambda(\alpha)$, and $M^{-T}\Z^2=\Lambda^*$: the right-hand side is $\frac{1}{V\alpha}\Theta_{\Lambda^*}(1/\alpha)$.
For $A$, Lemma~\ref{ch1:lem:autodual} gives $\Theta_{A^*}=\Theta_A$.
\end{humanproof}

\begin{uniremark}\tierL{} The same three steps prove the Poisson summation formula in $\R^d$: for a lattice $\Lambda\subset\R^d$ and a Schwartz
function $f$, $\sum_{x\in\Lambda}f(x)=\operatorname{covol}(\Lambda)^{-1}\sum_{y\in\Lambda^*}\hat f(y)$. With
$f(x)=e^{-\pi\alpha|x|^2}$, whose Fourier transform is $\alpha^{-d/2}e^{-\pi|\xi|^2/\alpha}$, it yields
$\Theta_\Lambda(\alpha)=\operatorname{covol}(\Lambda)^{-1}\alpha^{-d/2}\,\Theta_{\Lambda^*}(1/\alpha)$, the identity behind the Jacobi theta
transformation and the functional equation of Epstein zeta functions (standard; see also~\cite{montgomery} for the planar case).
\end{uniremark}

\begin{example}[Jacobi theta function]\label{ch1:ex:jacobi}
In dimension~1, the same proof gives, for $\vartheta(\alpha)=\sum_{n\in\Z}e^{-\pi\alpha n^2}$, the relation
$\vartheta(1/\alpha)=\sqrt\alpha\,\vartheta(\alpha)$. It shows that $\vartheta(\alpha)\sim\alpha^{-1/2}$ as $\alpha\to0$: for small $\alpha$ the sum
``sees'' the Gaussian integral. \tierC{}
\end{example}

\begin{numbox}[: the Poisson formula with its controls (code \texttt{code/ch1\_theta.py})]
For $A$, we compare $\Theta_A(\alpha)$ with $\alpha^{-1}\Theta_A(1/\alpha)$ (positive control), and then with $\Theta_A(1/\alpha)$ \emph{without} the factor
$\alpha^{-1}$ (negative control). Actual output of the script (in double precision):
\begin{center}\small
\begin{tabular}{@{}rlll@{}}
\toprule
$\alpha$ & $\Theta_A(\alpha)$ & gap to $\alpha^{-1}\Theta_A(1/\alpha)$ & gap without the factor \\
\midrule
$0.5$ & $2.008477918598$ & $4.4\cdot10^{-16}$ & $1.0$ \\
$2$ & $1.004238959299$ & $2.2\cdot10^{-16}$ & $1.0$ \\
$3$ & $1.000112671200$ & $8.9\cdot10^{-16}$ & $2.0$ \\
\bottomrule
\end{tabular}
\end{center}
The negative control is discriminating ($\alpha\neq1$); at $\alpha=1$ both gaps vanish, so that point tests nothing.
The script also compares $\Theta_A(\alpha)\le\Theta_{\Z^2}(\alpha)$ at eight values of $\alpha\in[0.1, 5]$ (always true), for example
$\Theta_A(1)=1.159595267$ against $\Theta_{\Z^2}(1)=1.180340599$. \tierX{}: this illustrates Montgomery's theorem below, without proving it.
\end{numbox}

\begin{theorem}[Montgomery, 1988]\label{ch1:thm:montgomery}
For every lattice $\Lambda$ in the plane of covolume $1$ and every $\alpha>0$, $\Theta_A(\alpha)\le\Theta_\Lambda(\alpha)$, with equality if and only if $\Lambda$ is
the image of $A$ under a rotation. \tierL{} (cited second-hand)~\cite{montgomery}
\end{theorem}

Montgomery thus solves the optimality problem \emph{among lattices}, for every Gaussian interaction. We return to this in §\ref{ch1:sec:univ}
to pass from Gaussians to completely monotone potentials. The difficulty lies elsewhere: comparing $A$ not with other lattices, but with
\emph{all} point configurations of density $1$.

\section{Completely monotone functions}
\label{ch1:sec:cm}

\subsection{Definition and first examples}

\begin{definition}[Completely monotone function]\label{ch1:def:cm}
A function $g:(0,\infty)\to\R$ is \emph{completely monotone} (CM) if it is of class $\mathcal C^\infty$ and
$(-1)^n g^{(n)}(t)\ge0$ for every $n\in\N$ and every $t>0$.
\end{definition}

The case $n=0$ says $g\ge0$; the case $n=1$ says that $g$ is nonincreasing; the case $n=2$ says that $g$ is convex; and so on.
The inequalities are \emph{non-strict}: constants are CM. In the \texttt{openai/math} repository, a function that is $\mathcal C^\infty$ on $(0,\infty)$,
nonnegative and CM is called \emph{admissible}~\cite{oai}; since $n=0$ gives nonnegativity, the two notions coincide. This coincidence is
itself a Lean theorem: \lean{ProgCM.isCompletelyMonotoneOn_iff_admissible} \tierA{} (library of programme P3~\cite{p3}).

\begin{example}[Exponential, powers]\label{ch1:ex:exemples}
\begin{enumerate}
\item $g(t)=e^{-at}$, $a\ge0$: $(-1)^ng^{(n)}(t)=a^ne^{-at}\ge0$. \tierC{}
\item $g(t)=t^{-s}$, $s\ge0$: $(-1)^ng^{(n)}(t)=s(s+1)\cdots(s+n-1)\,t^{-s-n}\ge0$ (immediate induction). \tierC{}
For $s=\frac12$ this is the Coulomb repulsion $1/r$ in the variable $t=r^2$. The CM property of $t\mapsto t^{-s/2}$ (Riesz) is also
checked by the kernel: \lean{TriangularRiesz.riesz_admissible} \tierA{}.
\end{enumerate}
\end{example}

\begin{example}[Counterexamples]\label{ch1:ex:contre}
$g(t)=t$ is not CM ($-g'=-1<0$), nor is $\cos t$ (it takes negative values), nor $e^{t}$ ($-g'<0$).
A subtler counterexample, of physical origin (a charge--quadrupole combination), is
$q(t)=t^{-1/2}-2(t+1)^{-1/2}+(t+4)^{-1/2}$: we have $q(3)=3^{-1/2}-2\cdot4^{-1/2}+7^{-1/2}\approx0.5774-1+0.3780<0$,
so $q$ is not nonnegative, let alone CM. The rigorous verification of $q(3)<0$ and the non-CM property of $q$ are checked by the kernel
(\lean{ProgCM.quadrupole_not_cm}) \tierA{}.
\end{example}

\subsection{Closure properties: human proofs}

\begin{proposition}[Sum, product]\label{ch1:prop:cloture}
Let $f,g$ be CM and $c\ge0$. Then $f+g$, $cf$ and $fg$ are CM. \tierC{} (Lean version: programme P3~\cite{p3}, \tierA{})
\end{proposition}

\begin{humanproof}
The sum and multiplication by a nonnegative constant are immediate by linearity of differentiation.
For the product, the Leibniz rule gives
\[
  (-1)^n(fg)^{(n)}=\sum_{k=0}^{n}\binom nk\ \bigl[(-1)^kf^{(k)}\bigr]\ \bigl[(-1)^{n-k}g^{(n-k)}\bigr],
\]
since $(-1)^n=(-1)^k(-1)^{n-k}$. Each bracket is nonnegative, hence so is the sum.
\end{humanproof}

A \emph{Bernstein function} is the right notion for composing with a CM function.

\begin{definition}[Bernstein function]\label{ch1:def:bernstein}
A $\mathcal C^\infty$ function $h:(0,\infty)\to[0,\infty)$ is a \emph{Bernstein function} if its derivative $h'$ is completely monotone
(Lean definition \lean{IsBernsteinOn}; Theorem~\ref{ch1:thm:comp} additionally requires $h>0$). In particular $h$ is nondecreasing.
\end{definition}

\begin{example}\label{ch1:ex:bernstein}
The following are Bernstein functions: $h(t)=at+b$ ($a,b\ge0$); $h(t)=t^\alpha$ for $0<\alpha\le1$, since
$h'(t)=\alpha\,t^{-(1-\alpha)}$ is CM by Example~\ref{ch1:ex:exemples}; $h(t)=1-e^{-at}$ ($a>0$), since $h'=ae^{-at}$; $h(t)=\ln(1+t)$,
since $h'(t)=1/(1+t)$. In particular $h(t)=\sqrt t$ (the case $\alpha=\frac12$) is one. \tierC{}
\end{example}

\begin{uniremark}\tierL{} Bernstein functions admit a L\'evy--Khintchine representation: $h:(0,\infty)\to[0,\infty)$ is a Bernstein function
if and only if $h(t)=a+bt+\int_{(0,\infty)}(1-e^{-st})\,\nu(ds)$ with $a,b\ge0$ and $\nu$ a positive measure satisfying
$\int_{(0,\infty)}\min(1,s)\,\nu(ds)<\infty$. This follows by integrating the Bernstein--Widder representation (Theorem~\ref{ch1:thm:bernstein})
of the CM function $h'$; for instance $\sqrt t=\frac{1}{2\sqrt\pi}\int_0^\infty(1-e^{-st})\,s^{-3/2}\,ds$ (standard).
\end{uniremark}

\begin{theorem}[Subordination: CM $\circ$ Bernstein]\label{ch1:thm:comp}
If $f$ is CM and $h$ is a \emph{strictly positive} Bernstein function, then $f\circ h$ is CM. \tierC{} Lean version:
\texttt{ProgCM.IsCompletelyMonotoneOn\allowbreak .comp\_bernstein} \tierA{}, whose proof follows exactly the induction below~\cite{p3}.
\end{theorem}

\begin{humanproof}
The composition $f\circ h$ is well defined (since $h>0$) and of class $\mathcal C^\infty$. We prove by strong induction on $n$ the property
\[
  \mathcal P(n) :\quad \text{for every CM } f \text{ and every Bernstein } h,\ \ (-1)^n(f\circ h)^{(n)}\ge0\ \text{on }(0,\infty).
\]
It is essential to quantify over \emph{all} functions $f$ and $h$: the induction hypothesis will be applied to a function other than $f$.

\emph{Base case.} $\mathcal P(0)$: $f\circ h\ge0$ since $f\ge0$.

\emph{Inductive step.} Assume $\mathcal P(k)$ for every $k\le n$. Let $f$ be CM and $h$ Bernstein. We have $(f\circ h)'=(f'\circ h)\,h'$, hence
\[
  -(f\circ h)'=G\cdot h',\qquad G=(-f')\circ h .
\]
Now $-f'$ is CM (it is $f$ differentiated once, with the sign: $(-1)^m(-f')^{(m)}=(-1)^{m+1}f^{(m+1)}\ge0$). By $\mathcal P(k)$ applied to the CM function $-f'$,
$(-1)^kG^{(k)}\ge0$ for every $k\le n$. On the other hand $h'$ is CM by definition. Differentiating $n$ times and using the Leibniz rule,
\[
  (-1)^{n+1}(f\circ h)^{(n+1)}=(-1)^{n}\bigl(G\,h'\bigr)^{(n)}=\sum_{k=0}^{n}\binom nk\ \bigl[(-1)^kG^{(k)}\bigr]\ \bigl[(-1)^{n-k}(h')^{(n-k)}\bigr]\ \ge0 .
\]
Hence $\mathcal P(n+1)$.
\end{humanproof}

\begin{corollary}[The square-root case]\label{ch1:cor:sqrt}
If $f$ is CM, then $t\mapsto f(\sqrt t)$ is CM. In particular, for $\kappa\ge0$, $t\mapsto e^{-\kappa\sqrt t}$ is CM, and so is the
planar Yukawa potential $e^{-\kappa r}/r$ written in the variable $t=r^2$: $g_\kappa(t)=e^{-\kappa\sqrt t}\,t^{-1/2}$. \tierC{}
(Lean: \lean{ProgCM.IsCompletelyMonotoneOn.comp_sqrt}, \lean{ProgCM.phi_cm} \tierA{}.)
\end{corollary}

\begin{humanproof}
Theorem~\ref{ch1:thm:comp} with $h=\sqrt{\ \cdot\ }$ (Example~\ref{ch1:ex:bernstein}) gives the first sentence. With $f(u)=e^{-\kappa u}$ we obtain that
$e^{-\kappa\sqrt t}$ is CM. The Yukawa potential is the \emph{product} of this function with $t^{-1/2}$, which is CM by Example~\ref{ch1:ex:exemples};
Proposition~\ref{ch1:prop:cloture} concludes.
\end{humanproof}

\begin{remark}[Direct check on one case]
For $\psi(t)=e^{-\sqrt t}$ we have $\psi'(t)=-\frac{1}{2\sqrt t}e^{-\sqrt t}$ and $\psi''(t)=\bigl(\frac1{4t}+\frac1{4t^{3/2}}\bigr)e^{-\sqrt t}$. The signs
$-\psi'\ge0$ and $\psi''\ge0$ are as announced. A machine checks the signs up to order $8$ at seven points
(script \texttt{code/ch1\_theta.py}): $e^{-\sqrt t}$ passes, $\cos t$ (negative control) fails at every order \tierX{}. This is not a proof:
the proof is the induction above.
\end{remark}

\begin{physbox}[: why physicists meet CM functions]
A potential that decays with distance while flattening out ``more and more slowly'' (Coulomb, screened Yukawa, dipole) is typically CM.
The five physical potentials of Table~\ref{ch1:tab:flux} include dipolar excitons $2\bigl(1/r-1/\sqrt{r^2+d^2}\bigr)$ (doi:\,10.5281/zenodo.23646103) and the Yukawa potential
(doi:\,10.5281/zenodo.23646105). In both cases admissibility is a \tierA{} theorem; the physical consequence (the triangular lattice has
minimal energy) remains conditional on \texttt{openai/math}.
\end{physbox}

\subsection{Bernstein's theorem and Gaussian mixtures}

\begin{theorem}[Bernstein--Widder]\label{ch1:thm:bernstein}
A function $g:(0,\infty)\to\R$ is completely monotone if and only if there exists a positive measure $\mu$ on $[0,\infty)$ such that
\[
  g(t)=\int_{[0,\infty)}e^{-\alpha t}\,d\mu(\alpha)\qquad(t>0).
\]
The measure $\mu$ is then unique. \tierL{} (classical statement, cited second-hand)~\cite{bernstein,widder}
\end{theorem}

The geometric reading is essential: with $t=|x|^2$, we have $g(|x|^2)=\int e^{-\alpha|x|^2}d\mu(\alpha)$, that is, \emph{every completely monotone potential
is a positive mixture of Gaussians}. The ``converse'' direction is easy and we prove it; the forward direction (CM $\Rightarrow$ mixture) is a deep
result, which we cite, and which is \emph{not} formalised in the work described here.

\begin{proposition}[Mixtures are CM (easy direction)]\label{ch1:prop:facile}
Let $\mu$ be a positive measure on $[0,\infty)$ such that $\int e^{-\alpha t}d\mu(\alpha)<\infty$ for every $t>0$. Then $g(t)=\int e^{-\alpha t}d\mu(\alpha)$ is CM.
\tierC{} (Lean: \lean{ProgCM.laplace_measure} \tierA{}, a statement of this same direction.)
\end{proposition}

\begin{humanproof}
Fix $t_0>0$. For $t\ge t_0$ and $n\in\N$, we have $0\le\alpha^ne^{-\alpha t}\le C_n\,e^{-\alpha t_0/2}$ with
$C_n=\sup_{\alpha\ge0}\alpha^ne^{-\alpha t_0/2}<\infty$, and $\alpha\mapsto e^{-\alpha t_0/2}$ is $\mu$-integrable by hypothesis. The theorem on differentiation under the
integral sign (uniform domination on $[t_0,\infty)$, a consequence of Lebesgue's dominated convergence theorem) shows, by induction on $n$, that $g$ is of class $\mathcal C^\infty$ and
$(-1)^ng^{(n)}(t)=\int\alpha^ne^{-\alpha t}d\mu(\alpha)\ge0$.
\end{humanproof}

\begin{example}[A CM potential that the closure rules do not see]\label{ch1:ex:dipole}
For $c>0$, $\varphi(t)=t^{-1/2}-(t+c)^{-1/2}$ is a \emph{difference} of two CM functions, and nothing in Proposition~\ref{ch1:prop:cloture}
allows us to conclude. But Euler's formula $u^{-1/2}=\frac{1}{\sqrt\pi}\int_0^\infty \alpha^{-1/2}e^{-\alpha u}d\alpha$ ($u>0$) gives
\[
  \varphi(t)=\frac{1}{\sqrt\pi}\int_0^\infty\alpha^{-1/2}\bigl(1-e^{-c\alpha}\bigr)\,e^{-\alpha t}\,d\alpha ,
\]
a mixture with \emph{positive} weight ($1-e^{-c\alpha}>0$): $\varphi$ is CM by Proposition~\ref{ch1:prop:facile}. This is the dipolar exciton potential,
up to the factor $2$. \tierC{}
\end{example}

\paragraph{From the energy of Gaussians to that of mixtures.} Let $g(t)=\int e^{-\pi\alpha t}d\nu(\alpha)$ be a mixture, with $\nu$ positive.
For a finite configuration $C_R$, by positivity (Tonelli's theorem) the double sum can be moved under the integral:
\[
  E_g(C;R)=\int E_{e^{-\pi\alpha\,\cdot}}(C;R)\,d\nu(\alpha).
\]
Fatou's lemma (valid in $[0,\infty]$) then gives $E_g(C)\ge\int\liminf_RE_{e^{-\pi\alpha\cdot}}(C;R)\,d\nu(\alpha)=\int E_{e^{-\pi\alpha\cdot}}(C)\,d\nu(\alpha)$.
If we know that, for each $\alpha$, no configuration has a Gaussian energy lower than that of $A$, we deduce
$E_g(C)\ge\int(\Theta_A(\alpha)-1)\,d\nu(\alpha)=E_g(A)$, by Tonelli on the lattice sum. \tierC{}
\emph{The Gaussian case is therefore the heart of the problem}; it is this case, for every value of $\alpha$, that the theorem of the next part asserts.
The same idea, applied to \emph{lattices} with Montgomery's Theorem~\ref{ch1:thm:montgomery}, gives: $E_g(A)\le E_g(\Lambda)$ for every lattice $\Lambda$ of
covolume $1$ and every CM $g$ (a mixture $\nu$ of Gaussians), since $E_g(\Lambda)=\int(\Theta_\Lambda(\alpha)-1)\,d\nu(\alpha)$. \tierC{}\ (conditional on~\ref{ch1:thm:montgomery} and~\ref{ch1:thm:bernstein}, \tierL{})

\section{Universal optimality}
\label{ch1:sec:univ}

\subsection{Statement}

Cohn and Kumar~\cite{ck} defined a notion of \emph{universal optimality}: a configuration is universally optimal if it minimises the energy for
\emph{every} completely monotone potential of the squared distance, simultaneously. They conjectured that this happens for many remarkable
configurations, including the triangular lattice in the plane. Cohn, Kumar, Miller, Radchenko and Viazovska~\cite{ckmrv} proved the conjecture for the $E_8$ lattice in dimension $8$ and the Leech lattice
in dimension $24$ \tierL{}, building on Viazovska's work~\cite{viazovska}. Dimension $2$ remained open in the literature we have read.

The \texttt{openai/math} repository contains a statement and a proof in Lean~4 of the planar case~\cite{oai}. Here is the statement as it appears in the
challenge file, first in English and then in Lean.

\begin{theorem}[Universal optimality of the triangular lattice (\texttt{openai/math})]\label{ch1:thm:univ}
Let $g$ be admissible (Definition~\ref{ch1:def:cm}) and $C\subset\R^2$ a locally finite configuration of centred density $1$. Then
\[
  \mathcal E_g(A)\ \le\ E_g(C)\qquad\text{and}\qquad \mathcal E_g(A)=E_g(A),\qquad\text{where}\quad \mathcal E_g(A)=\sum_{x\in A\setminus 0}g(|x|^2),
\]
with values in $[0,\infty]$, and $E_g(\cdot)$ the lower energy over closed centred disks (Lean: \lean{latticeEnergy g} for $\mathcal E_g(A)$ and \lean{energy g} for $E_g$). \tierA{} \emph{for the Lean statement}
\lean{OAI.AtomicTriangular.universal_energy_minimum}, \textbf{conditional}: the proof is AI-generated and has not been peer reviewed.
\end{theorem}

\begin{leanbox}[: the main statement (extract from the challenge file, where the proof is \lean{sorry} by construction)]
\begin{lstlisting}[language=Lean]
theorem universal_energy_minimum (g : ℝ → ℝ) (C : Set Plane)
    (hg : AdmissiblePotential g) (hC : LocallyFinite C) (hd : DensityOne C) :
    latticeEnergy g ≤ energy g C ∧ latticeEnergy g = energy g A
\end{lstlisting}
\end{leanbox}

Three reading remarks, essential for what follows.
\begin{enumerate}
\item The energy is a \emph{limit inferior} over \emph{closed disks centred at $0$}. The CKMRV definition~\cite{ckmrv} additionally requires the limit to exist for the
optimal configuration. The notion used here is therefore weaker \tierC{}.
\item The theorem does not say that $A$ is the only minimiser, and uniqueness is false in general: corollary D3 of textbook v1~\cite{v1} establishes
non-uniqueness by removing one point of $A$ (a ``vacancy''), which changes neither the density nor the lower energy. \tierA{} (conditional)
\item For $s\le2$ (Riesz), the energy of $A$ is infinite (Proposition~\ref{ch1:prop:riesz}): the inequality is then true but void. Content appears for $s>2$.
\end{enumerate}

\subsection{The thread of the formalisations}

The rest of the book follows a natural sequence: each step takes the previous result as a black box.
Table~\ref{ch1:tab:flux} summarises it. The DOI numbers and tiers are those of the record documents.

\begin{table}[ht]
\centering\small
\begin{tabular}{@{}p{0.04\textwidth}p{0.50\textwidth}p{0.17\textwidth}p{0.20\textwidth}@{}}
\toprule
\# & Step & Tier & Reference \\
\midrule
1 & Upstream theorem (\texttt{openai/math}), planar case & \tierL{} and conditional \tierA{} & \cite{oai} \\
2 & Corollaries D1 to D4: Riesz, every density, CKMRV sense and non-uniqueness, open or closed disks & \tierA{}, conditional on 1 & \cite{v1} (textbook v1) \\
3 & Five physical potentials (dipolar excitons, Yukawa, Rytova--Keldysh, $\sigma=3$ enstrophy, Bogoliubov structure factor) & \tierA{} for the replayed modules (dipolar, Yukawa, Keldysh), \tierAs{} for $\sigma=3$ and Bogoliubov; Keldysh$\leftrightarrow$Struve: \tierC{} & 5 Zenodo deposits \\
4 & CM/Bernstein library (P3): subordination, equivalence with \lean{AdmissiblePotential} & \tierA{} & \cite{p3} \\
5 & Quantum surprise: the zero-point energy, QU-1 false for 8 CM potentials out of 12; certificate at $a=1$; no triangular$\to$oblique transition & \tierX{}, then \tierAs{} (logic) and \tierC{} (ball-arithmetic integrals); negative result & Zenodo \\
6 & Sandier--Serfaty (Abrikosov lattices): audit, partial port to Lean, failure of the extension to the 3D case & \tierC{}, partial \tierAs{}; negative results & Zenodo \\
7 & Navier--Stokes (corpus \texttt{openai/NavierStokesAndEuler}): rebuild and kernel replay & \tierA{} for the kernel replay of the corpus statements; faithfulness to Fefferman's alternative (C): \tierC{}; no ``Clay problem solved'' & doi:\newline 10.5281/\allowbreak zenodo.23661381 \\
8 & Population genetics: Moran fixation, Price equation, isothermal theorem & \tierA{} & Zenodo \\
\bottomrule
\end{tabular}
\caption{The sequence of formalisations, from the foundations (this chapter) to the applications.}
\label{ch1:tab:flux}
\end{table}

\begin{warnbox}[: what the table does not claim]
No row means ``conjecture proved'' beyond its tier. Rows 5 and 6 contain \emph{negative} results: they
say that certain plausible extensions fail (or are not established), and this is information of the same value as the positive results.
Step~7 is an audit of a third-party corpus: the credit belongs to its authors. We do not claim that a Clay problem is solved.
\end{warnbox}

\section{Lean 4 for the student}
\label{ch1:sec:lean}

\subsection{Propositions and proofs}

Lean~4~\cite{lean4} is both a programming language and a \emph{proof assistant}. For an undergraduate reader,
the following dictionary is enough to read the statements of this book (this is the Curry--Howard correspondence).
\begin{itemize}
\item A \emph{proposition} is a mathematical statement; in Lean, it is an object of type \lean{Prop}.
\item A \emph{proof} of the proposition $P$ is an object \lean{h : P}. The Lean kernel checks that a term really has the announced type.
\item $P\wedge Q$ is written \lean{P ∧ Q}: a proof is a \emph{pair} $\langle$\lean{hP}, \lean{hQ}$\rangle$.
\item $P\Rightarrow Q$ is written \lean{P → Q}: a proof is a \emph{function} that turns every proof of $P$ into a proof of $Q$.
\item $\forall n,\ P(n)$ is written \lean{∀ n, P n}: a proof is a function that assigns to each $n$ a proof of $P(n)$.
\end{itemize}
The Mathlib library~\cite{mathlib} provides many already proved theorems (analysis, algebra).

\subsection{A minimal example, with its actual output}

The following file was compiled as is with Lean 4.34.0-rc2, without any library (output reproduced identically by the independent auditor). It contains three true theorems and two false
``theorems'' $1=2$ that can be proved only by cheating: one with \lean{sorry}, the other with an axiom added by the user.

\begin{leanbox}[: \texttt{T.lean} (compiled on 2026-10-26)]
\begin{lstlisting}[language=Lean]
theorem and_swap' (p q : Prop) (h : p ∧ q) : q ∧ p := ⟨h.2, h.1⟩

theorem add_comm' (a b : Nat) : a + b = b + a := Nat.add_comm a b

theorem two_mul' (n : Nat) : 2 * n = n + n := by
  omega

theorem bad : 1 = 2 := by sorry

axiom cheat : False
theorem bad2 : 1 = 2 := cheat.elim

#print axioms and_swap'
#print axioms two_mul'
#print axioms bad
#print axioms bad2
\end{lstlisting}
\end{leanbox}

The first two theorems are \emph{proof terms}: the object is written directly. The third uses a \emph{tactic} (\lean{by omega}):
a small program that \emph{constructs} the proof term; the kernel then checks the term produced. Here is the actual output:

\begin{lstlisting}[language={}]
T.lean:8:8: warning: declaration uses `sorry`
'and_swap'' does not depend on any axioms
'two_mul'' depends on axioms: [propext, Quot.sound]
'bad' depends on axioms: [sorryAx]
'bad2' depends on axioms: [cheat]
\end{lstlisting}

The compiler's exit code was $0$: \textbf{Lean accepts a file that contains \lean{sorry}}, emitting only a warning.
\emph{A zero exit code does not mean ``everything is proved''.}

\subsection{Axioms, and the gate}

Every proof in mathematics rests on axioms. The classical mathematics of Mathlib uses three: \lean{propext} (two equivalent
propositions are equal), \lean{Classical.choice} (the axiom of choice) and \lean{Quot.sound} (compatibility of quotients). We call them the \emph{whitelist}:
\[
  \{\texttt{propext},\ \texttt{Classical.choice},\ \texttt{Quot.sound}\}.
\]
Any other axiom in the footprint of a theorem disqualifies it. In the output above, the footprint of \lean{bad} contains \lean{sorryAx} (the internal axiom
hidden behind \lean{sorry}) and that of \lean{bad2} contains the user axiom \lean{cheat}: both false theorems are therefore \emph{detected}.
This is the \emph{negative control} of the axiom gate; the \emph{positive control} is that \lean{and_swap'} and \lean{two_mul'} pass. A measurement without these two controls proves nothing.

\begin{warnbox}[: \texttt{\#print axioms} is not enough]
The A2 audit~\cite{a2} of the \texttt{openai/NavierStokesAndEuler} corpus (2\,440 modules rebuilt, kernel replay of $53\,752+34\,648$ declarations) showed that a declaration
\emph{added directly to the kernel with its type checking disabled} (option \lean{debug.skipKernelTC}) is displayed by \texttt{\#print axioms} as
``does not depend on any axioms'', even though it can establish \lean{False}. The audit's positive control (such a false declaration, deliberately fabricated) was
\emph{invisible} to \texttt{\#print axioms} and \emph{caught} by the full kernel replay (\lean{leanchecker}) (a \tierX{}-tier observation in the A2 paper: it concerns a deliberately invalid declaration). The lesson: the axiom audit is necessary,
not sufficient.
\end{warnbox}

The \tierA{} gate of this book therefore consists of \emph{four} controls (compilation, then the three checks of the preface):
\begin{enumerate}
\item compilation of the project without errors;
\item the axiom list of each retained theorem is contained in the whitelist;
\item full \emph{kernel replay} of the declarations by the tool \lean{leanchecker}, which rechecks every term without trusting the compiler;
\item a textual search for the words \lean{skipKernelTC}, \lean{addDecl}, \lean{native_decide} and \lean{implemented_by} in the sources, which are the known backdoors.
\end{enumerate}
To this is added an organisational rule: whoever \emph{produces} a proof is not the one who \emph{verifies} it (here, two separate AI agents, and never a human).

When the \lean{leanchecker} replay has not been done, we write \tierAs{}: axiom list, planted witnesses and independent recompilation, without kernel replay.

\subsection{Reading the upstream statement; a short corollary}

Once the main theorem is available, a corollary often reduces to checking a hypothesis. Here is the Riesz case,
quoted from the file \lean{formal_cert/Triangular/Riesz.lean} of the repository (the ellipsis marks an omitted part):

\begin{leanbox}[: Riesz corollary (extract, layout modified)]
\begin{lstlisting}[language=Lean]
noncomputable def riesz (s : ℝ) : ℝ → ℝ := fun t => t ^ (-(s / 2))

theorem riesz_admissible {s : ℝ} (hs : 0 < s) : AdmissiblePotential (riesz s) := by
  ...

theorem triangular_riesz_optimal {s : ℝ} (hs : 0 < s) (C : Set Plane)
    (hC : LocallyFinite C) (hd : DensityOne C) :
    latticeEnergy (riesz s) ≤ energy (riesz s) C :=
  (universal_energy_minimum (riesz s) C (riesz_admissible hs) hC hd).1
\end{lstlisting}
\end{leanbox}

The last line ``applies'' the main theorem to $g=\text{riesz}\ s$ by supplying the three proofs of the hypotheses, and then takes the first component
\lean{.1} of the conjunction. This is exactly a \emph{composition of functions} in the sense of the Curry--Howard dictionary.
The hypothesis \lean{riesz_admissible} is the complete monotonicity property of Example~\ref{ch1:ex:exemples} (2). The ``every density'' corollary (D2) is re-exported in Chapter~5 under the name \lean{ALIX}.

\subsection{The induction of §\ref{ch1:sec:cm}, on the Lean side}

The subordination Theorem~\ref{ch1:thm:comp} has been formalised~\cite{p3}. The general statement and its square-root special case are the following
(extracts from the file \lean{Bernstein.lean} in the folder \lean{formal_cert/Programmes/}\allowbreak\lean{CompletelyMonotone/}, proofs omitted):

\begin{leanbox}[: subordination (statements)]
\begin{lstlisting}[language=Lean]
theorem IsCompletelyMonotoneOn.comp_bernstein {f g : ℝ → ℝ} (hf : IsCompletelyMonotoneOn f)
    (hg : IsBernsteinOn g) (hpos : ∀ t, 0 < t → 0 < g t) :
    IsCompletelyMonotoneOn (fun t => f (g t))

theorem IsCompletelyMonotoneOn.comp_sqrt {f : ℝ → ℝ} (hf : IsCompletelyMonotoneOn f) :
    IsCompletelyMonotoneOn (fun t => f (Real.sqrt t))
\end{lstlisting}
\end{leanbox}

One recognises the structure of the human theorem. The hypotheses \lean{hf}, \lean{hg} and \lean{hpos} are the three
hypotheses of the theorem ($f$ CM, $h$ Bernstein, $h>0$), \lean{fun t => f (g t)} is the function $f\circ h$, and \lean{comp_sqrt} is obtained from the general case by
supplying the proof that $\sqrt{\ \cdot\ }$ is a Bernstein function and strictly positive on $(0,\infty)$. The repository states that the proof proceeds by strong
induction on the order, quantified over all CM functions $f$, with the formula $(f\circ h)'=-((-f')\circ h)\cdot h'$ and the Leibniz lemma: this is
the same reasoning as the human proof. \tierA{} for the statements above, \tierC{} for our human write-up, whose correspondence with the
Lean proof is a comment by the author, not a theorem.

\begin{remark}[Why \lean{hpos} appears in Lean]
In paper mathematics, one tacitly assumes that $f\circ h$ makes sense. Lean does not: $f$ is a function defined everywhere on $\R$, but
constrained only on $(0,\infty)$. If $h\equiv0$ (which is a Bernstein function), $f\circ h$ would be the spurious constant $f(0)$, about which the hypothesis says nothing. The hypothesis
$h>0$ is not nitpicking: it is what makes the statement \emph{true}.
\end{remark}

\begin{warnbox}[: summary of the tiers of this chapter]
\begin{itemize}
\item \tierL{}: Poisson summation on the torus (Fourier's theorem), Montgomery's theorem and the Bernstein--Widder theorem (forward direction), both cited second-hand, the CKMRV result in dimensions $8$ and $24$.
\item \tierC{}: all human proofs of this chapter (counting, Riesz energy, Poisson summation, closure of CM functions, subordination, mixtures, Fatou reduction).
\item \tierA{} (with kernel replay for the Triangular modules, the CM/Bernstein library and P7; see the following chapters): the Lean statements cited (equivalence with \lean{AdmissiblePotential}, subordination, \lean{laplace_measure}, \lean{quadrupole_not_cm}, the upstream statement), the last one being \emph{conditional}.
\item \tierX{}: the script \texttt{code/ch1\_theta.py}, which illustrates without proving.
\end{itemize}
\end{warnbox}

\subsection{The Elenchus tiers}
\label{ch1:sec:tiers}

The Elenchus method~\cite{elenchus} attaches to each statement of the book one of four tiers, according to the \emph{nature of the available proof}.

\begin{center}\small
\begin{tabular}{@{}p{0.09\textwidth}p{0.40\textwidth}p{0.38\textwidth}@{}}
\toprule
Tier & What is guaranteed & What is not \\
\midrule
\tierA{} & The Lean statement is accepted by the kernel, with the four controls of the gate above & That the Lean statement faithfully expresses the intended mathematical or physical statement \\
\tierL{} & The result appears in a source that was read and precisely referenced (for example a theorem of a published paper) & That we have rechecked it ourselves \\
\tierC{} & A complete human argument, written in this book, reviewed by an independent AI agent & Any machine verification, and any human review \\
\tierX{} & A floating-point computation, with positive and negative controls & Nothing: a numerical exploration never supports a claim \\
\bottomrule
\end{tabular}
\end{center}

Two composition rules, used throughout what follows.
\begin{enumerate}
\item \emph{The tier of a consequence is the weakest of the tiers of its hypotheses.} A corollary of \texttt{openai/math} is \emph{conditional} \tierA{}: the kernel has
checked the implication, not the upstream hypothesis, whose proof has not been peer reviewed.
\item \emph{A negative result is a result.} When a plausible extension fails, we write it down with its tier.
\end{enumerate}

\begin{warnbox}[: adequacy of the statement]
The kernel certifies that the statement \lean{universal_energy_minimum}, \emph{with the repository's definitions}, follows from the whitelisted axioms. It does not certify
that \lean{DensityOne} is the density the physicist has in mind, nor that \lean{energy} (a limit inferior over closed centred disks) is the CKMRV energy. No human
has audited this adequacy. This is why every chapter repeats the tier, and it is the reason for the word ``conditional''.
\end{warnbox}

\section*{Exercises}
\addcontentsline{toc}{section}{Exercises}
\label{ch1:sec:exos}

Hints follow the list of exercises. The exercises are of increasing difficulty (from * to ***).

\begin{exercise}[*]\label{ch1:ex:covol}
Show that the change-of-basis matrix between two bases of the same lattice lies in $\mathrm{GL}_2(\Z)$ and deduce that the covolume does not depend on the basis.
Compute the covolume and the minimal distance of the lattice $\Z u\oplus\Z v$ with $u=(1,0)$, $v=(1/2,\sqrt3/2)$, then normalise it to density $1$. Recover $a^2=2/\sqrt3$.
\end{exercise}

\begin{exercise}[*]\label{ch1:ex:voisins}
In $A$, determine all vectors of norm $a\sqrt3$ and their number (hint: $m^2+mn+n^2=3$). Compare with the case of $\Z^2$ (shell $|x|^2=2$). Justify that $a>1$: at equal covolume, the minimal distance of $A$ exceeds that of $\Z^2$.
\end{exercise}

\begin{exercise}[**]\label{ch1:ex:comptage}
Write out the proof of Proposition~\ref{ch1:prop:comptage} for $\Lambda=\Z^2$ and deduce that $\#\{(m,n)\in\Z^2 : m^2+n^2\le R^2\}=\pi R^2+O(R)$.
Is this error $O(R)$ optimal? (Look up the literature on the \emph{Gauss circle problem}.)
\end{exercise}

\begin{exercise}[**]\label{ch1:ex:epstein}
Prove that $\sum_{x\in A\setminus0}|x|^{-s}$ diverges for $s\le2$ and converges for $s>2$. Using $t^{-s/2}=\frac{\pi^{s/2}}{\Gamma(s/2)}\int_0^\infty\alpha^{s/2-1}e^{-\pi\alpha t}d\alpha$,
express this sum as an integral of $\Theta_A(\alpha)-1$ against $\alpha^{s/2-1}d\alpha$ and recover the condition $s>2$ from the behaviour of $\Theta_A$ at $0$ and at $+\infty$.
\end{exercise}

\begin{exercise}[**]\label{ch1:ex:poisson}
(a) Prove the relation $\vartheta(1/\alpha)=\sqrt\alpha\,\vartheta(\alpha)$ of Example~\ref{ch1:ex:jacobi} by redoing the one-dimensional proof.
(b) Deduce an algorithm that computes $\vartheta(\alpha)$ to within $10^{-12}$ by summing only over $|n|\le 6$, whatever $\alpha>0$. (c) Compute
$\Theta_{\Z^2}(1)=\vartheta(1)^2$ and compare with the value $1.1803\ldots$ in the text.
\end{exercise}

\begin{exercise}[*]\label{ch1:ex:clotures}
Show that if $f$ is CM and $a>0$, $b\ge0$, then $t\mapsto f(at+b)$ is CM. Deduce that $t\mapsto(t+1)^{-s}e^{-a t}$ is CM for $s,a\ge0$.
Show by a computation that $t\mapsto t^{-1}e^{-t}$ is CM as well, without using the composition theorem.
\end{exercise}

\begin{exercise}[**]\label{ch1:ex:comp}
(a) Check that $h(t)=\ln(1+t)$ is a Bernstein function. (b) Deduce that $t\mapsto(\ln(1+t))^{-s}$ is CM for $s\ge0$ (pay attention to the positivity of $h$ on $(0,\infty)$).
(c) Show that $t\mapsto1/(1+\sqrt t)$ is CM by applying Corollary~\ref{ch1:cor:sqrt}.
(d) Why can the hypothesis $h>0$ not be omitted in Theorem~\ref{ch1:thm:comp}? (Consider $h\equiv0$.)
\end{exercise}

\begin{exercise}[***]\label{ch1:ex:melange}
Let $g(t)=t^{-1/2}-2(t+1)^{-1/2}+(t+4)^{-1/2}$. Noting that $g$ is a mixture of Gaussians with a \emph{signed weight}, explain why positivity of the weight fails,
and recover in this way that $g$ is not CM (compare with $g(3)<0$).
\end{exercise}

\begin{exercise}[*]\label{ch1:ex:lean}
Among the following four outputs of \texttt{\#print axioms}, which satisfy the axiom gate? Which satisfy the \emph{full} four-control gate?
(i) \texttt{[propext, Classical.choice, Quot.sound]}; (ii) \texttt{[propext, sorryAx]}; (iii) \texttt{does not depend on any axioms}; (iv) \texttt{[cheat, propext]}.
Explain why (iii) alone is not enough to conclude.
\end{exercise}

\begin{exercise}[**]\label{ch1:ex:fatou}
Let $\Lambda$ be a lattice and $g\ge0$ an arbitrary potential. Show that the disk energy $E_g(\Lambda;R)$ defined above converges, as $R\to\infty$,
to $\sum_{a\in\Lambda\setminus0}g(|a|^2)$ (finite or not). Thus, for a lattice, the lower energy and the lattice energy coincide.
\end{exercise}

\begin{exercise}[(*) University exercise: Poisson summation in $\R^d$]
Let $\Lambda=M\Z^d\subset\R^d$ be a lattice, $M\in\mathrm{GL}_d(\R)$, with dual $\Lambda^*=M^{-T}\Z^d$ and $\Theta_\Lambda(\alpha)=\sum_{x\in\Lambda}e^{-\pi\alpha|x|^2}$.
(a) Adapting the three steps of the proof of Theorem~\ref{ch1:thm:poisson}, prove that for every $\alpha>0$
\[
  \Theta_\Lambda(\alpha)=\frac{1}{\operatorname{covol}(\Lambda)}\,\alpha^{-d/2}\,\Theta_{\Lambda^*}\!\Bigl(\frac1\alpha\Bigr),\qquad \operatorname{covol}(\Lambda)=|\det M|.
\]
(b) Deduce that $\Theta_\Lambda(\alpha)\sim\operatorname{covol}(\Lambda)^{-1}\alpha^{-d/2}$ as $\alpha\to0^+$, and recover the case $d=1$ of Example~\ref{ch1:ex:jacobi}.
(c) Show that if $\Lambda^*$ is the image of $\Lambda$ under an orthogonal map (so that $\operatorname{covol}(\Lambda)=1$), then $\alpha\mapsto\alpha^{d/4}\Theta_\Lambda(\alpha)$ is invariant under $\alpha\mapsto1/\alpha$.
\end{exercise}

\begin{exercise}[(*) University exercise: Exponentials of Bernstein functions]
Let $h$ be a Bernstein function with $h>0$ on $(0,\infty)$.
(a) Show that $t\mapsto e^{-u\,h(t)}$ is CM for every $u\ge0$, and that $t\mapsto 1/h(t)$ is CM.
(b) Show that if $h_1,h_2$ are Bernstein functions with $h_2>0$, then $h_1\circ h_2$ is a Bernstein function. (Hint: $(h_1\circ h_2)'=(h_1'\circ h_2)\,h_2'$; use Theorem~\ref{ch1:thm:comp} and Proposition~\ref{ch1:prop:cloture}.)
(c) Deduce that $t\mapsto\bigl(\ln(1+\sqrt t)\bigr)^{-1}$ and $t\mapsto e^{-\sqrt{t+t^{1/2}}}$ are CM.
\end{exercise}

\subsection*{Hints}

\begin{itemize}
\item[\ref{ch1:ex:covol}.] If $M,M'$ are two basis matrices, $M'=MP$ with $P$ integral and $M=M'P'$ with $P'$ integral: $PP'=I$, so $\det P\det P'=1$ with two
integers, hence $\det P=\pm1$. For the last part: $\det(u,v)=\sqrt3/2$; rescale by the factor $a$ with $a^2\sqrt3/2=1$.
\item[\ref{ch1:ex:voisins}.] $m^2+mn+n^2=3$: $(1,1)$, $(-1,-1)$, $(2,-1)$, $(-2,1)$, $(1,-2)$, $(-1,2)$, i.e.\ six vectors. Since $a^2=2/\sqrt3>1$ ($4>3$), $a>1$; the
minimal distance of $\Z^2$ is $1$.\item[\ref{ch1:ex:comptage}.] Take $\delta=\sqrt2$ for $\Z^2$: $\pi(R-\sqrt2)^2\le N(R)\le\pi(R+\sqrt2)^2$, hence $N(R)=\pi R^2+O(R)$; optimality of the exponent $1$ is not asked for (the best known exponent is below $1$, see the literature).
\item[\ref{ch1:ex:epstein}.] For divergence, reuse the lower bound on $\#S_k$. For the integral: $\Theta_A(\alpha)-1\sim6e^{-\pi a^2\alpha}$ as $\alpha\to\infty$ (convergence)
and $\Theta_A(\alpha)\sim\alpha^{-1}$ as $\alpha\to0$ (by Poisson); the integral $\int_0\alpha^{s/2-1}\alpha^{-1}d\alpha$ converges iff $s/2-2>-1$, i.e.\ $s>2$.
\item[\ref{ch1:ex:poisson}.] (b) For $\alpha\ge1$ sum directly; for $\alpha<1$ apply $\vartheta(\alpha)=\alpha^{-1/2}\vartheta(1/\alpha)$.
(c) $\vartheta(1)=1.0864348\ldots$, hence $\Theta_{\Z^2}(1)=1.18034\ldots$
\item[\ref{ch1:ex:clotures}.] $f(at+b)$: differentiate $n$ times and factor out $a^n$. For the last one: product of two CM functions (Proposition~\ref{ch1:prop:cloture}); or a direct computation via Leibniz.
\item[\ref{ch1:ex:comp}.] (b) Set $f(u)=u^{-s}$, which is CM. (d) For $h\equiv 0$, $f\circ h$ is not defined on $(0,\infty)$ (or takes a spurious value $f(0)$ if one extends).
\item[\ref{ch1:ex:melange}.] $t^{-1/2}=\pi^{-1/2}\int\alpha^{-1/2}e^{-\alpha t}d\alpha$; the weight of $g$ is $\pi^{-1/2}\alpha^{-1/2}(1-2e^{-\alpha}+e^{-4\alpha})$, which is negative for
some $\alpha$ (study $1-2x+x^4$ with $x=e^{-\alpha}$). Injectivity of the Laplace transform (including for signed measures) shows that $g$ has no representation with a positive weight; more directly, $g(3)<0$ concludes.
\item[\ref{ch1:ex:lean}.] (i) passes. (ii) fails (\lean{sorryAx}). (iii) passes the axiom gate alone, but the declaration might have been inserted by bypassing the
kernel's type checker. (iv) fails.
\item[\ref{ch1:ex:fatou}.] The term of $x\in\Lambda_R$ is $\sum_{a\in(\Lambda_R-x)\setminus0}g(|a|^2)$: it is bounded above by $\sum_{a\ne0}g(|a|^2)$, whence the upper bound on $E_g(\Lambda;R)$.
For the lower bound, fix $r>0$: if $|x|\le R-r$, the term of $x$ contains all $a\ne0$ with $|a|\le r$. Hence
$E_g(\Lambda;R)\ge\frac{N_\Lambda(R-r)}{N_\Lambda(R)}\sum_{0<|a|\le r}g(|a|^2)$; the ratio tends to $1$ (Proposition~\ref{ch1:prop:comptage}). Let $R$ and then $r$ tend to infinity.
\end{itemize}
